\newpage
\subsubsection{Функции переходов и выходов}
\label{sec:functions}

В таблице~\ref{tab:functions} указаны пары $\delta(q,x)/\lambda(q,x)$ для нажатий кнопок. Строка задает текущее состояние $q$, столбец~--- входной символ $x$.

\begin{table}[H]
	\centering
	\caption{Функции переходов и выходов при нажатии кнопок}
	\label{tab:functions}
	\small\setstretch{1.12}
	\renewcommand{\arraystretch}{1.1}
	\setlength{\tabcolsep}{9pt}
	\begin{tabular}{|c|cccc|}
		\hline
		$q \backslash x$ & $a$ & $b$ & $c$ & $d$ \\
		\hline
		$S_{0}$ & $S_{1}/z_{0}$ & $S_{4}/z_{0}$ & $S_{7}/z_{0}$ & $S_{0}/z_{0}$ \\
		$S_{1}$ & $S_{2}/z_{0}$ & $S_{1}/z_{1}$ & $S_{2}/z_{0}$ & $S_{2}/z_{0}$ \\
		$S_{2}$ & $S_{3}/z_{0}$ & $S_{2}/z_{2}$ & $S_{3}/z_{0}$ & $S_{3}/z_{0}$ \\
		$S_{3}$ & $S_{10}/z_{0}$ & $S_{3}/z_{3}$ & $S_{10}/z_{0}$ & $S_{10}/z_{0}$ \\
		$S_{4}$ & $S_{5}/z_{4}$ & $S_{6}/z_{5}$ & $S_{0}/z_{0}$ & $S_{0}/z_{0}$ \\
		$S_{5}$ & $S_{4}/z_{4}$ & $S_{4}/z_{4}$ & $S_{4}/z_{4}$ & $S_{4}/z_{4}$ \\
		$S_{6}$ & $S_{4}/z_{5}$ & $S_{4}/z_{5}$ & $S_{6}/z_{6}$ & $S_{4}/z_{5}$ \\
		$S_{7}$ & $S_{7}/z_{7}$ & $S_{8}/z_{0}$ & $S_{11}/z_{10}$ & $S_{0}/z_{0}$ \\
		$S_{8}$ & $S_{8}/z_{8}$ & $S_{9}/z_{0}$ & $S_{9}/z_{0}$ & $S_{9}/z_{0}$ \\
		$S_{9}$ & $S_{9}/z_{9}$ & $S_{7}/z_{0}$ & $S_{7}/z_{0}$ & $S_{7}/z_{0}$ \\
		$S_{10}$ & $S_{0}/z_{0}$ & $S_{0}/z_{0}$ & $S_{0}/z_{0}$ & $S_{0}/z_{0}$ \\
		$S_{11}$ & $S_{11}/z_{11}$ & $S_{7}/z_{0}$ & $S_{7}/z_{0}$ & $S_{7}/z_{0}$ \\
		\hline
	\end{tabular}
\end{table}

Для событий времени и такта без события правила одинаковы во всех состояниях $q\in Q$:
\[
	\begin{array}{c|ccc}
		x & \tau & \alpha & \varnothing \\ \hline
		\delta(q,x) & q & q & q \\
		\lambda(q,x) & z_{12} & z_{13} & z_0
	\end{array}
\]

Каждой паре из $Q\times X$ соответствует ровно одно следующее состояние и один выход. Все состояния достижимы из $S_0$ по нажатию кнопок, а из каждого режима можно вернуться в $S_0$.

Для просмотра текущего времени при работающем секундомере из $S_0$ последовательно нажимаются $b$, $a$, $b$, $c$. Автомат проходит состояния $S_4$, $S_5$, $S_4$ и $S_0$. На последнем переходе выдается $z_0$, поэтому секундомер продолжает отсчет. Чтобы снова увидеть показание секундомера и остановить его, нужно дважды нажать $b$.

\newpage
В таблице~\ref{tab:binary-functions} те же переходы по кнопкам записаны в двоичном виде. В каждой ячейке сначала приведен код следующего состояния, а затем код выхода. Для входов $100$, $101$ и $110$ код состояния сохраняется, а код выхода равен соответственно $1100$, $1101$ и $0000$.

\begin{table}[H]
	\centering
	\caption{Двоичные коды переходов и выходов при нажатии кнопок}
	\label{tab:binary-functions}
	\small\setstretch{1.12}
	\renewcommand{\arraystretch}{1.1}
	\begin{tabularx}{\textwidth}{|c|*{4}{>{\centering\arraybackslash}X}|}
		\hline
		Состояние & $000$ & $001$ & $010$ & $011$ \\
		\hline
		$S_{0}$ & $0001/0000$ & $0100/0000$ & $0111/0000$ & $0000/0000$ \\
		$S_{1}$ & $0010/0000$ & $0001/0001$ & $0010/0000$ & $0010/0000$ \\
		$S_{2}$ & $0011/0000$ & $0010/0010$ & $0011/0000$ & $0011/0000$ \\
		$S_{3}$ & $1010/0000$ & $0011/0011$ & $1010/0000$ & $1010/0000$ \\
		$S_{4}$ & $0101/0100$ & $0110/0101$ & $0000/0000$ & $0000/0000$ \\
		$S_{5}$ & $0100/0100$ & $0100/0100$ & $0100/0100$ & $0100/0100$ \\
		$S_{6}$ & $0100/0101$ & $0100/0101$ & $0110/0110$ & $0100/0101$ \\
		$S_{7}$ & $0111/0111$ & $1000/0000$ & $1011/1010$ & $0000/0000$ \\
		$S_{8}$ & $1000/1000$ & $1001/0000$ & $1001/0000$ & $1001/0000$ \\
		$S_{9}$ & $1001/1001$ & $0111/0000$ & $0111/0000$ & $0111/0000$ \\
		$S_{10}$ & $0000/0000$ & $0000/0000$ & $0000/0000$ & $0000/0000$ \\
		$S_{11}$ & $1011/1011$ & $0111/0000$ & $0111/0000$ & $0111/0000$ \\
		\hline
	\end{tabularx}
\end{table}
